\section{Introduction}
Les chapitres précédents ont permis de poser l'ensemble des bases de cette recherche : le chapitre 1 a présenté l'état de l'art, 
le chapitre 2 a présenté le protocole expérimental, le chapitre 3 a détaillé la mise en œuvre concrète de ce protocole sur le terrain, tandis que le chapitre 4 a présenté la modélisation et la conception de l'application de pilotage de l'irrigation et de la fertilisation.

Ce chapitre  rassemble et discute l'ensemble des résultats obtenus au cours de l'expérimentation, qu'ils soient issus du suivi du sol, 
de la croissance des cultures, du comportement phytosanitaire des plots, du rendement final, ou du fonctionnement de l'application en question. 
Il permet une interprétation honnête et nuancée de l'effet réel de l'amendement du sol et du pilotage de l'irrigation et de la fertilisation sur les performances agronomiques des trois cultures étudiées : le gombo, le navet et la laitue.

Le chapitre se répartit  en cinq grandes sections. En premier lieu, les résultats relatifs à l'amendement du sol sont présentés (section 5.2) ; ensuite, viennent les résultats des courbes d'évolution des paramètres étudiés (section 5.3) ; 
puis, en section 5.4, nous avons les résultats phytosanitaires ; enfin, les résultats des rendements sont aussi présentés en section 5.5 avant une conclusion partielle (section 5.6) menant aux recommandations et perspectives de ce travail.
